\section{Related Work}
\label{sec:rl_related_work}

This section reviews prior work related to the adaptation methods developed in this chapter. The discussion is organized around five themes: hierarchical and temporal abstraction in reinforcement learning, hybrid planning and learning systems, creative problem solving and action discovery, transfer learning and continual adaptation, and prior approaches to novelty handling in open-world settings.

\paragraph{Hierarchical reinforcement learning and temporal abstraction.}
The executor formalism used in this chapter draws on the options framework~\citep{sutton1999options}, which extends reinforcement learning to include temporally extended actions defined by initiation sets, policies, and termination conditions. Options enable agents to reason and learn at multiple levels of temporal granularity, and have been shown to improve exploration and transfer in hierarchical settings. Hierarchical deep RL methods~\citep{kulkarni2016hierarchical} decompose tasks into subgoal sequences, improving sample efficiency in complex domains. The connection between temporal abstraction and symbolic planning has been formalized by \citet{konidaris2018skills}, who show that the symbols necessary and sufficient for planning can be derived from the initiation sets and effect distributions of a set of skills, and by \citet{konidaris2016constructing}, who develops the skill-symbol loop in which new skills give rise to new symbolic representations. The executor discovery problem introduced in this chapter can be viewed as the inverse of the skill-symbol loop: given a symbolic operator whose executor has failed, the agent must learn a new skill (executor) that satisfies the operator's specification.

\paragraph{Hybrid planning and learning.}
Integrating symbolic planning with learned execution has been explored in several contexts. Task and motion planning (TAMP) systems~\citep{garrett2021integrated,kaelbling2011hierarchical} combine symbolic task planning with continuous motion planning, but typically assume complete and correct domain models. PDDLStream~\citep{garrett2020pddlstream} extends this paradigm by integrating symbolic planners with blackbox samplers, while \citet{silver2021learning} learn symbolic operators for task and motion planning from demonstrations. \citet{konidaris2014constructing} show how to construct symbolic representations suitable for high-level planning from continuous robotic experience, establishing theoretical connections between skill learning and symbolic abstractions that underpin the operator--executor framework used in this chapter. Plan monitoring~\citep{fritz2007monitoring} and plan repair~\citep{fox2006plan} address execution failures within the symbolic planning paradigm, either by detecting when a plan becomes suboptimal or by modifying a plan to account for unexpected states. However, these approaches operate entirely at the symbolic level and cannot acquire new low-level behaviors when the failure stems from a mismatch between the domain model and the environment dynamics. The SPOTTER framework~\citep{sarathy2021spotter} takes a step toward integrating planning with learning by using reinforcement learning to extend symbolic planning operators, and MacGyver problems~\citep{sarathy2018macgyver} highlight the broader challenge of resourceful problem solving that motivates novelty-aware architectures. More recent hybrid approaches leverage approximate symbolic models for skill diversity~\citep{guan2022leveraging}, learn symbolic options to facilitate deep RL~\citep{jin2021creativity}, or integrate relational planning with reinforcement learning~\citep{kokel2021reprel}. \citet{lyu2019sdrl} demonstrate interpretable deep RL that leverages symbolic planning for both structure and data efficiency. \citet{illanes2020symbolic} use symbolic plans as high-level instructions for RL agents, providing hierarchical structure that improves learning. \citet{yang2018peorl} integrate symbolic planning with hierarchical RL for robust decision-making, and \citet{thierauf2024fixing} combine RL with symbolic planning in object-based action spaces for plan repair. These approaches demonstrate the complementary strengths of planning and learning, but they generally assume a fixed domain model and are not evaluated in open-world settings where plan execution failures are the norm. RAPid-Learn differs from these approaches in that it explicitly targets the executor discovery problem: when an operator fails due to novelty, it instantiates a focused reinforcement learning problem guided by the symbolic structure of the failed plan.

\paragraph{Creative problem solving and action discovery.}
A closely related line of work addresses the problem of discovering new actions or behaviors when an agent's existing repertoire proves insufficient. \citet{gizzi2019creative} introduce a framework for creative problem solving by robots using action primitive discovery, in which novel actions are discovered through segmenting the low-level trajectory representations of known symbolic actions. This work establishes the connection between symbolic action representations and subsymbolic trajectory analysis that is also central to the operator--executor decomposition used in this chapter. \citet{gizzi2021behavior} extend this approach through behavior babbling, in which agents discover new actions by exploring variations of known behaviors, analogous to how motor babbling enables infants to discover new movement capabilities. These discovered actions are represented both symbolically and subsymbolically, enabling the agent to re-plan in novel scenarios with expanded knowledge.

\citet{gizzi2022cps} provide a comprehensive survey and framework for creative problem solving (CPS) in artificially intelligent agents, situating action discovery within a broader landscape that includes analogical reasoning, conceptual blending, and constraint relaxation. They identify the key challenge that motivates the present chapter: when an agent's existing knowledge is insufficient to solve a task due to environmental change, the agent must generate novel solutions that go beyond recombining known elements. \citet{gizzi2022lifelong} further extend this to a lifelong learning setting, using world models to improve performance in novelty resolution across successive encounters. The adaptation methods developed in this chapter address a specific instance of creative problem solving---recovery from execution failures---using reinforcement learning to discover new executors guided by the symbolic structure of the failed plan.

\paragraph{Transfer learning and continual adaptation.}
When an agent encounters a changed environment, one natural approach is to transfer knowledge from the original task to the new setting. \citet{taylor2009transfer} provide a comprehensive survey of transfer learning for RL domains. Policy reuse~\citep{glatt2017policy} transfers a learned policy and selects between the old and new policies during execution. Actor-critic transfer methods~\citep{clegg2017haptics} transfer both the policy and value function networks. Progressive neural networks~\citep{rusu2016progressive} avoid catastrophic forgetting by freezing previously learned columns and adding new capacity for new tasks, while elastic weight consolidation~\citep{kirkpatrick2017overcoming} protects important weights from modification. Curriculum learning~\citep{narvekar2020curriculum,narvekar2016source} provides a complementary approach by structuring the sequence of training tasks to improve sample efficiency and transfer, with direct connections to how knowledge can be accumulated across novelty encounters. Meta-RL approaches~\citep{yu2020metaworld} attempt to adapt to multiple MDPs and solve unseen tasks drawn from the same distribution, but assume gradual variation rather than abrupt novelty. \citet{khetarpal2022continual} provide a comprehensive review of continual RL methods, noting that most approaches assume continuous evolution and struggle with sudden, structural changes. \citet{lorang2022speeding} propose methods for speeding up continual learning through information gains in novel experiences, demonstrating that prior knowledge about novelty structure can accelerate adaptation. These methods address the problem of leveraging prior experience, but they treat adaptation as a monolithic learning problem over the full state-action space. When novelty affects only a specific component of the task, this global approach is inefficient: the agent must relearn or readjust its entire policy rather than targeting the specific behavior that failed. The experimental results in \cref{sec:rl_eval_rapid} confirm this limitation, showing that transfer learning baselines require one to two orders of magnitude more experience than RAPid-Learn to achieve comparable post-novelty performance.

\paragraph{Open-world novelty handling.}
The DARPA SAIL-ON program~\citep{senator2019science} motivated the development of agent architectures for detecting and accommodating novelty in open-world environments. \citet{boult2021towards} and \citet{langley2020open} provide formal frameworks for reasoning about novelty, while \citet{muhammad2021novelty} present an agent architecture that combines novelty detection with reactive accommodation strategies. \citet{balloch2023neuro} develop neuro-symbolic world models that adapt internal representations in response to detected novelty. \citet{thai2022architecture} present an architecture for novelty handling in multi-agent stochastic environments, evaluated in open-world Monopoly. Complementary agent-modeling work learns interpretable disposition models of other agents online, hypothesizing latent internal states when observable features alone cannot explain behavior---an important capability when novelty is behavioral rather than only environmental~\citep{lymperopoulos2024agentmodels}. More recently, \citet{lorang2024adapting} demonstrate the utility of hybrid hierarchical RL and symbolic planning for open-world adaptation, and \citet{lorang2025curiosity} propose curiosity-driven imagination for discovering plan operators and learning associated policies in open-world settings. The cognitive architecture presented in \cref{sec:architecture} extends this line of work by integrating symbolic and neural novelty detection with a cascading recovery pipeline that progresses from lightweight symbolic strategies to learning-based executor discovery, and by evaluating on a substantially larger and more diverse set of novelty scenarios than prior work.

\paragraph{Gap addressed by this chapter.}
Across these lines of work, a common limitation is the absence of a principled mechanism for converting a global task failure into a targeted learning problem. Symbolic methods can detect and localize failures but cannot synthesize new behaviors. Learning methods can acquire new behaviors but lack the structural guidance needed for efficient adaptation. Action discovery methods~\citep{gizzi2019creative,gizzi2021behavior} demonstrate that novel behaviors can be discovered through exploration, but they do not leverage the symbolic structure of a failed plan to focus the search. The RAPid-Learn framework addresses this gap by leveraging the action abstractions developed in \cref{sec:action_abstractions}---the operator--executor decomposition---to bridge symbolic failure diagnosis with focused reinforcement learning, enabling adaptation that is both targeted and sample-efficient.